\chapter{Implement functions from signatures}
\source{Sample Q2, p. 2 and Q4(b), p. 4; GenericMonads lecture.}
Write code directly in these answer blocks. Implement the described behavior,
not merely an expression that type-checks. No \code{undefined}, \code{error},
or intentional nontermination. The optional Haskell companion repeats these tasks.

\task{3.1. Collect present values}
Return all values in order when every element is \code{Just}; otherwise return
\code{Nothing}. An empty list succeeds with an empty result. Use explicit
recursion, not \code{sequence}, \code{traverse}, \code{mapM}, or equivalent helpers.
\par\noindent\begin{minipage}{\linewidth}
\begin{codeanswer}
collectMaybe :: [Maybe a] -> Maybe [a]
collectMaybe = undefined  -- YOUR CODE
\end{codeanswer}
\end{minipage}\par
\answer{[Types of the head, recursive result, and constructed payload.]}

\task{3.2. Generalize the operation}
Implement left-to-right traversal and action collection. Use explicit recursion
and do/bind/return; do not call traversal, sequencing, or monadic fold helpers.
For \code{gather}, do not delegate to \code{walkM}.
\par\noindent\begin{minipage}{\linewidth}
\begin{codeanswer}
walkM :: Monad m => (a -> m b) -> [a] -> m [b]
walkM = undefined  -- YOUR CODE

gather :: Monad m => [m a] -> m [a]
gather = undefined  -- YOUR CODE
\end{codeanswer}
\end{minipage}\par

\task{3.3. Rewrite using a library combinator}
Reimplement \code{gather} with exactly one call to library \code{mapM}, with no
explicit recursion. \code{mapM :: Monad m => (x -> m y) -> [x] -> m [y]}.
\par\noindent\begin{minipage}{\linewidth}
\begin{codeanswer}
gatherViaMapM :: Monad m => [m a] -> m [a]
gatherViaMapM = undefined  -- YOUR CODE
\end{codeanswer}
\end{minipage}\par
\answer{[What are x and y in this use of mapM? What type must its function argument have?]}

\clearpage
\task{3.4. More generic monadic operations}
Use only do/bind/return, ordinary application, and supplied arguments.
\code{composeM} runs the first action-producing function and feeds its result
to the second. \code{flattenM} removes one monadic layer. \code{liftTwoM}
performs its two actions left-to-right and combines their results.
\par\noindent\begin{minipage}{\linewidth}
\begin{codeanswer}
composeM :: Monad m => (a -> m b) -> (b -> m c) -> a -> m c
composeM = undefined  -- YOUR CODE

flattenM :: Monad m => m (m a) -> m a
flattenM = undefined  -- YOUR CODE

liftTwoM :: Monad m => (a -> b -> c) -> m a -> m b -> m c
liftTwoM = undefined  -- YOUR CODE
\end{codeanswer}
\end{minipage}\par

\task{3.5. A limit of polymorphism}
Is there a total implementation of \code{forall m a. Monad m => m a} for
every \code{m} and \code{a}? Explain using both \code{Maybe} and
\code{State ()}. Why does the existence of \code{Nothing} not settle the question?
\answer{[Reasoning about all monads and all result types.]}

\task{3.6. Changed implementation}
Implement \code{repeatAction}: collect the results of performing one action
\code{n} times, in order. For \code{n <= 0}, perform no actions. Use explicit
recursion; no \code{replicateM}, \code{mapM}, \code{sequence}, or \code{traverse}.
\par\noindent\begin{minipage}{\linewidth}
\begin{codeanswer}
repeatAction :: Monad m => Int -> m a -> m [a]
repeatAction = undefined  -- YOUR CODE
\end{codeanswer}
\end{minipage}\par
\answer{[Predict repeatAction 2 "ab", including order, before running it.]}
